\chapter{PHỤ LỤC}
\label{ch:phu-luc}

\section{Hướng dẫn gán nhãn cảm xúc trong miền dữ liệu}

Tập trong miền dữ liệu được gán theo ba lớp NEGATIVE, NEUTRAL và POSITIVE. Nhãn phản ánh tác
động kỳ vọng của nội dung đối với doanh nghiệp hoặc mã cổ phiếu liên quan, không chỉ
phản ánh sắc thái của một câu riêng lẻ. Người gán nhãn phải đọc tiêu đề và phần nội
dung có sẵn, ghi nhận trường hợp có nhiều sự kiện hoặc hướng tác động không rõ ràng,
đồng thời không suy luận nhãn từ biến động giá xảy ra sau thời điểm bài viết được
đăng.

Các mẫu được mô hình gán nhãn sơ bộ rồi được một người rà soát toàn bộ. Báo cáo cuối
phải ghi số người rà soát, số mẫu hợp lệ, phân bố ba lớp, số mẫu phải chỉnh sửa và
nguyên tắc xử lý mẫu khó. Do chỉ có một người rà soát, không tính Cohen's kappa hoặc
Fleiss' kappa.

\section{Hồ sơ thiết kế hệ thống}

\subsection{Tác nhân và chức năng}

Hệ thống có một tác nhân chính là người nghiên cứu. Người nghiên cứu cung cấp dữ liệu,
khóa cấu hình và kiểm tra các báo cáo sinh ra từ quy trình. Các chức năng chính gồm:

\begin{itemize}[nosep]
  \item Thu thập và kiểm tra dữ liệu tin tức, dữ liệu giá;
  \item Chuẩn hóa văn bản, loại bản ghi trùng và lưu nguồn gốc và khả năng truy xuất;
  \item Gán tin với mã cổ phiếu và phiên giao dịch theo mốc cắt;
  \item Suy luận xác suất ba lớp cảm xúc bằng PhoBERT;
  \item Tạo đặc trưng theo mã và ngày, huấn luyện mô hình dự báo;
  \item Đánh giá theo thứ tự thời gian và xuất bản kê, chỉ số, biểu đồ.
\end{itemize}

\subsection{Sơ đồ luồng xử lý}

Sơ đồ dưới đây mô tả luồng xử lý của hệ thống. Khối đánh giá chỉ nhận dữ liệu dự báo
ngoài mẫu và bản kê của cùng một lần chạy; không dùng dữ liệu tập giữ lại để điều chỉnh
ngược cấu hình.

\begin{figure}[H]
  \centering
  \begin{tikzpicture}[
      node distance=5mm,
      block/.style={
        rectangle, rounded corners=2pt, draw=black!60, thick,
        text width=10cm, minimum height=9mm, align=center,
        fill=gray!8, inner sep=4pt},
      decision/.style={
        rectangle, rounded corners=2pt, draw=blue!65, very thick,
        text width=10cm, minimum height=9mm, align=center,
        fill=blue!8, inner sep=4pt},
      arr/.style={-{Stealth[length=2.5mm]}, thick, black!65}
    ]
    \node[block] (input) {Tin tức và dữ liệu giá};
    \node[block, below=of input] (quality) {Kiểm tra lược đồ, mốc thời gian, mã và bản ghi trùng};
    \node[decision, below=of quality] (features) {Tạo đặc trưng cảm xúc và đặc trưng thị trường theo mã, ngày};
    \node[decision, below=of features] (models) {Huấn luyện đối chứng cơ sở, mô hình chỉ giá và LSTM hai nhánh};
    \node[block, below=of models] (evaluation) {Cuốn chiếu theo thời gian, tập giữ lại, loại bỏ đặc trưng và phân tích sai số};
    \node[block, below=of evaluation] (report) {Chỉ số, biểu đồ, bản kê và báo cáo tái lập};
    \draw[arr] (input) -- (quality);
    \draw[arr] (quality) -- (features);
    \draw[arr] (features) -- (models);
    \draw[arr] (models) -- (evaluation);
    \draw[arr] (evaluation) -- (report);
  \end{tikzpicture}
  \caption{Sơ đồ luồng xử lý và báo cáo kết quả của hệ thống.}
  \label{fig:phu-luc-luong-xu-ly}
\end{figure}
\noindent\textit{Nguồn: Tác giả xây dựng.}

\subsection{Các mô-đun chính}

\begin{description}[style=nextline, nosep]
  \item[\texttt{stf.data.news}] Đọc và chuẩn hóa dữ liệu tin, xử lý mốc thời gian, tiêu đề,
  nội dung và URL.
  \item[\texttt{stf.data.prices}] Đọc chuỗi giá OHLCV theo mã cổ phiếu và kiểm tra độ
  phủ lịch giao dịch.
  \item[\texttt{stf.sentiment.dataset}] Đọc tập nhãn, chuẩn hóa nhãn về ba lớp và tạo
  các tập huấn luyện, xác thực, kiểm thử theo thời gian khi có cột ngày.
  \item[\texttt{stf.sentiment.model}] Tinh chỉnh PhoBERT, lưu điểm kiểm và sinh xác
  suất cảm xúc cho ba lớp.
  \item[\texttt{stf.sentiment.metrics}] Tính F1 vĩ mô cố định trên ba lớp và chỉ số theo lớp.
\end{description}

\section{Cấu hình tái lập}

Snapshot dữ liệu đầu vào gồm \texttt{data/raw/news/articles.parquet} cho tin tức,
\texttt{data/raw/news/listings.parquet} cho liên kết tin với mã, các tệp giá theo mã
trong \texttt{data/raw/prices/}, cùng dữ liệu nhãn trong \texttt{data/labeled/}.
Mỗi lần chạy cần lưu phiên bản mã nguồn, phiên bản thư viện, hạt giống, điểm kiểm, mốc cắt,
phạm vi ngày, danh sách mã, mã băm dữ liệu và các tham số huấn luyện. Bản kê là nguồn
đối chiếu chính cho mọi con số xuất hiện trong phần kết quả thực nghiệm.

Tiền xử lý chỉ được ước lượng trên phần tập huấn luyện của từng cửa sổ. Tập giữ lại cuối chỉ được đọc ở
bước báo cáo. Nếu thay đổi nguồn dữ liệu, mốc cắt, ngưỡng tạo nhãn hoặc cách chia tập,
phải tạo bản kê mới và không gộp kết quả giữa hai cấu hình.

\section{Cài đặt và chạy hệ thống}

Môi trường được quản lý bằng \texttt{uv} và yêu cầu Python từ 3.12. Các lệnh kiểm tra
offline tối thiểu:

\begin{verbatim}
uv sync
uv run python -m stf.cli verify
uv run python -m stf.cli sentiment-smoke --n 60
\end{verbatim}

Huấn luyện PhoBERT cần một tệp nhãn đã hoàn tất với cột \texttt{text} và
\texttt{label}, nên chỉ được chạy sau khi khóa hướng dẫn gán nhãn và lưu bản kê
tương ứng. Bước huấn luyện chính thức nên thực hiện trên GPU; tên GPU, thời gian chạy
và điểm kiểm phải được lưu cùng kết quả.

\section{Mã nguồn}

Mã nguồn hệ thống được tổ chức trong gói \texttt{stf} của kho mã nguồn đề tài.
Các mô-đun thu thập dữ liệu, xử lý cảm xúc, đánh giá và giao diện dòng lệnh được đặt trong thư mục
\texttt{src/stf/}. Notebook huấn luyện GPU và hướng dẫn chạy được lưu trong thư mục
\texttt{notebooks/}. Phụ lục này mô tả giao diện và luồng xử lý; kết quả định lượng
chỉ được xác nhận bởi các sản phẩm đầu ra sinh từ một lần chạy có bản kê.
